\section{Introducción}
La presencia actual de materia bariónica y las propiedades conocidas de las partículas plantean una pregunta retrospectiva: ¿podemos utilizar nuestras masas de reposo, inventarios y cálculos, junto con las abundancias cosmológicas actuales, para reconstruir cuántos bariones y antibariones había durante la época de aniquilación y contrastar esa reconstrucción con lo que se conoce? El fondo cósmico de microondas y la nucleosíntesis primordial permiten inferir el pequeño exceso de número bariónico que permanece \citep{Fields2025,Planck2020}. Ese exceso y la población total de pares son magnitudes diferentes.

Un gas de resonancias hadrónicas (HRG, por sus siglas en inglés) utiliza las masas de reposo, degeneraciones y cargas de las especies de partícula para sumar sus poblaciones y contribuciones termodinámicas \citep{Vovchenko2019}. Esas masas ya formaban parte de nuestros cálculos. En este trabajo, las \emph{firmas de las partículas} se entienden como esas propiedades identificadoras del catálogo; no como señales espectrales originales medidas por nosotros. El modelo proporciona un vínculo matemático entre tales propiedades, una temperatura elegida y las poblaciones de equilibrio. Su utilidad a baja temperatura ha sido contrastada con cálculos de cromodinámica cuántica en red, con alcance dependiente del observable y de la proximidad a la transición \citep{Bazavov2012}.

La aniquilación elimina bariones y antibariones en igual número y conserva el exceso preexistente cuando no se viola el número bariónico. Las reacciones inversas también intervienen en la evolución térmica. \citet{Satarov2013} estudiaron ambas contribuciones y la expansión, y señalaron la importancia de las resonancias antibariónicas pesadas por encima de aproximadamente 100 MeV. Por ello, una comparación útil debe separar tres cuestiones: la reconstrucción de las poblaciones de equilibrio, la información que aportan las abundancias actuales y el historial de reacciones que exige un cálculo cinético.

\textbf{El objetivo principal fue comprobar si las propiedades de las especies de partícula y la abundancia bariónica actual permiten reconstruir las cantidades de bariones y antibariones en condiciones de la época de aniquilación, y contrastar esa reconstrucción con los resultados y aproximaciones publicados.} Para responderlo se resolvió una condición inversa que vinculó el exceso cosmológico observado con el potencial químico bariónico del HRG a ocho temperaturas. Se examinó además si la abundancia actual, por sí sola, permite recuperar una población inicial única. Como controles complementarios se compararon tres inventarios públicos, la aproximación nucleónica de la literatura y una referencia anterior de entropía \citep{Garcia2026}. La hipótesis de trabajo fue que una reconstrucción requiere declarar la temperatura y las condiciones térmicas, además de las masas y la abundancia presente.
